\documentclass[12pt]{scrreprt}
\usepackage{graphicx}
\usepackage{hyperref}
\usepackage{longtable}
\usepackage{array}
\usepackage{geometry}
\geometry{a4paper, margin=1in}

\title{\textbf{Software Requirement Specification (SRS)}\\
\large Software Component Cataloguing System}
\author{Prepared by: Shubham Singh}
\date{\today}

\begin{document}

\maketitle
\tableofcontents
\newpage

%--------------------------------------------------
\chapter{Introduction}

\section{Purpose}
The purpose of this Software Requirement Specification (SRS) document is to provide a detailed description of the \textbf{Software Component Cataloguing System}. This system is designed to store, manage, and retrieve reusable software components efficiently.

In modern software development, reuse of components significantly reduces development time and cost. However, without proper organization, locating reusable components becomes difficult. This system addresses that issue by providing a centralized catalogue.

\section{Scope}
The system will:
\begin{itemize}
    \item Maintain a repository of reusable software components
    \item Support both design (UML, ERD) and code components
    \item Allow keyword-based searching
    \item Track usage statistics of components
    \item Support hierarchical classification of components
\end{itemize}

\section{Intended Audience}
This document is intended for:
\begin{itemize}
    \item Developers
    \item System Designers
    \item Project Managers
    \item Testers
\end{itemize}

\vfill{}
\section{Definitions}
\begin{longtable}{|p{4cm}|p{10cm}|}
\hline
\textbf{Term} & \textbf{Description} \\
\hline
Component & A reusable piece of software (code or design) \\
Catalogue & Database storing components \\
Keyword & Tag used for searching components \\
Metadata & Information describing components \\
\hline
\end{longtable}

%--------------------------------------------------
\chapter{Overall Description}

\section{Product Perspective}
The Software Component Cataloguing System is a standalone application that acts as a repository for reusable components. It can also be integrated into larger development environments.

\section{Product Functions}
The major functions of the system include:
\begin{itemize}
    \item Adding new components
    \item Deleting existing components
    \item Updating component information
    \item Searching components using keywords
    \item Viewing usage statistics
\end{itemize}

\section{User Classes}
\begin{itemize}
    \item \textbf{Admin (Cataloguer):} Responsible for managing components
    \item \textbf{User (Developer):} Searches and uses components
\end{itemize}

\section{Operating Environment}
\begin{itemize}
    \item OS: Windows, Linux, MacOS
    \item Backend: Any server-side technology (Node.js / .NET / Java)
    \item Database: SQL / NoSQL
\end{itemize}

%--------------------------------------------------
\chapter{System Features}

\section{Feature 1: Add Component}
\textbf{Description:} Admin can add new components into the catalogue.

\textbf{Inputs:}
\begin{itemize}
    \item Component Name
    \item Type (Code/Design)
    \item Keywords
\end{itemize}

\textbf{Processing:}
\begin{itemize}
    \item Validate inputs
    \item Store in database
\end{itemize}

\textbf{Output:}
\begin{itemize}
    \item Confirmation message
\end{itemize}

\section{Feature 2: Search Component}
Users can search components using keywords.

\begin{itemize}
    \item Supports multiple keyword search
    \item Returns ranked results
\end{itemize}

\section{Feature 3: Usage Tracking}
System maintains:
\begin{itemize}
    \item Number of times component used
    \item Number of times searched but not used
\end{itemize}

%--------------------------------------------------
\chapter{Data Design}

\section{Database Schema}

\begin{longtable}{|p{3cm}|p{3cm}|p{7cm}|}
\hline
\textbf{Field} & \textbf{Type} & \textbf{Description} \\
\hline
ComponentID & Integer & Unique ID \\
Name & String & Component name \\
Type & String & Code/Design \\
Keywords & String & Search tags \\
UsageCount & Integer & Usage frequency \\
SearchMiss & Integer & Unused searches \\
\hline
\end{longtable}

%--------------------------------------------------
\chapter{System Models}

\section{Use Case Diagram}

\begin{figure}[h!]
\centering
\fbox{\parbox{12cm}{
Admin → Add/Delete Component \\
User → Search Component
}}
\caption{Use Case Diagram}
\end{figure}

\section{Data Flow Diagram}

\begin{figure}[h!]
\centering
\fbox{\parbox{12cm}{
User → Query → System → Database → Result
}}
\caption{DFD}
\end{figure}

\section{Architecture Diagram}

\begin{figure}[h!]
\centering
\fbox{\parbox{12cm}{
Frontend → Backend → Database
}}
\caption{System Architecture}
\end{figure}

%--------------------------------------------------
\chapter{Non-Functional Requirements}

\section{Performance}
\begin{itemize}
    \item Response time < 2 seconds
    \item Efficient search algorithm
\end{itemize}

\section{Security}
\begin{itemize}
    \item Authentication required
    \item Role-based access control
\end{itemize}

\section{Reliability}
\begin{itemize}
    \item 99\% uptime
\end{itemize}

\section{Scalability}
System should handle thousands of components efficiently.

%--------------------------------------------------
\chapter{Charts and Analysis}

\section{Component Usage Distribution}

\begin{figure}[h!]
\centering
\fbox{\parbox{10cm}{
Bar Chart showing usage frequency of components
}}
\caption{Usage Chart}
\end{figure}

%--------------------------------------------------
\chapter{Conclusion}

The Software Component Cataloguing System enhances software development by promoting reuse, reducing redundancy, and improving productivity. It provides an efficient mechanism for storing and retrieving reusable components.

\end{document}